\documentstyle[%


righttag%


,11pt%


,amssymb%


,russian%               remove in Eng.


%,izvru%                 Set izven% in Eng.


,izvru-ks%             for brief communications


]{amsart}





%





\newcommand{\const}{\operatorname{const}}


\newcommand{\re}{\operatorname{Re}}


\newcommand{\tts}[1]{}





\theoremstyle{definition}


\theorembodyfont{\rm}


\newtheorem{cornonum}{\hskip\parindent Следствие}


\renewcommand{\thecornonum}{}





%\hoffset=-15mm%                  For Eng.


\setcounter{page}{78}


\god{2002}


\nomer{6~(481)} %                        Remove in Eng.


%\nomer{Vol.\,46, No.\,6}%               Set in Eng.


%\str{\pageref{begin}--\pageref{end}}%  Set in Eng.


\udk{517.968}





\author{\it З.Б.\,Цалюк, С.С.\,Чистоклетова}


% NB:   ^^ remove in Eng.


\title{Асимптотика решений уравнений\\


Вольтерра--Гаммерштейна}





\date{}


%\pagestyle{myheadings}%         Set in Eng.


\pagestyle{plain} %              Remove in Eng.


\begin{document}


%\thispagestyle{plain}%          Set in Eng.


\maketitle


\label{begin}








Асимптотика решений линейных интегральных уравнений с разностным ядром


изучена сравнительно хорошо (напр., [1]--[8]). Здесь будет рассмотрено


уравнение


\begin{equation}


x(t)=\int_0^t K(t-s) [ x(s)+\varphi(x(s))] ds+f(t).


\end{equation}


Очевидно, в таком виде может быть записано любое уравнение


Вольтерра--Гаммерштейна с разностным ядром.


Ниже предполагается, что $K \in L_1


[0,\infty)$, $\varphi \in C(R)$, $f \in C[0,\infty)$. Как и в линейном


уравнении, асимптотика решений (1) (при заданной асимптотике $f$) связана с


нулями $1-\wh{K}(z)$, где $\wh{K}(z)= \int\limits_0^\infty e^{-zt} K(t)


dt$ --- преобразование Лапласа ядра $K(t)$.





\begin{thm}


Пусть $1-\wh{K}(z)$ не имеет нулей в полуплоскости $\re z \ge 0$ и


$\varphi (x)=o(x)$ при $|x| \to \infty$. Тогда при любой ограниченной $f(t)$


решение $(1)$ ограничено.


\end{thm}





Действительно, решение $x$ уравнения (1) является решением уравнения


\begin{equation}


x(t)=\int_0^t K(t-s) [ 1+q_\gamma (s) ] x(s)\,ds + \psi (t),


\end{equation}


где $q_\gamma (s) = \frac{\varphi[x(s)]}{x(s)}$, если $|x(s)| > \gamma$, и


$q_\gamma (s)=0$, если $|x(s)| \le \gamma$, а $\psi$ ---


некоторая ограниченная функция. По условию $K$ ---


устойчивое ядро ([9], гл.\,4, \S\,18, с.\,94), а $q_\gamma (s) \to 0$


равномерно по $s$ при $\gamma \to \infty$, значит, при достаточно


больших $\gamma$ ядро уравнения (2) устойчиво, а потому $x(t)$ ---


ограниченная функция.





\begin{thm}


Пусть $1-\wh{K}(z)$ имеет в полуплоскости $\re z\ge0$ конечное число нулей


$\lambda_1,\dotsc, \lambda_k$ и $\re\lambda_j>0$. Пусть далее $\varphi(x)=


O (|x|^\beta)$, $0 \le \beta < 1$, причем $\re\lambda_j \ne


\beta\mu$, где $\mu=\max \re\lambda_j$. Тогда при любой ограниченной $f$


решение


$$


x(t)=\sum_{\re\lambda_j > \beta \mu} P_j (t) e^{\lambda_j t} + O


(t^{\beta (n-1)} e^{\beta \mu t}),


$$


где $P_j (t)$ --- некоторые многочлены степени меньшей кратности корня


$\lambda_j$, а $n$ --- максимальная кратность тех $\lambda_j$,


для которых $\re\lambda_j = \mu$.


\end{thm}





Доказательство основано на следующих соображениях. В силу теоремы 1 


из [5] резольвента $R$ ядра $K$ имеет вид


\begin{equation}\label{r}


R(t)=\sum_{j=1}^k P_j (t) e^{\lambda_j t} + R_0 (t),


\ \ \text{где}\ \ R_0 \in L_1[0,\infty).


\end{equation}


Тогда


\begin{equation}\label{2}


x(t)=\int_0^t R(t-s) \varphi [ x(s) ] ds+ \psi(t),


\end{equation}


где $\psi(t)=f(t)+\int\limits_0^t R(t-s)f(s)ds=\sum\limits_j P_j (t) e^{


\lambda_j t} + u(t)$, $u(t)$ --- ограниченная функция.


По условию $|\varphi(x)|\le C_1[|x|^\beta+1]$, значит,


$$


|x(t)|\le C_1\int_0^t |R(t-s)|[|x^\beta(s)|+1]ds+|\psi(t)|.


$$


Функция $z(t)=C(t+1)^{n-1} e^{\mu t}$ при достаточно больших $C$


удовлетворяет неравенству


$$


z(t) \ge C_1\int_0^t |R(t-s)|[z^\beta(s)+1]ds+ |\psi(t)|.


$$


Отсюда и из теоремы об интегральном неравенстве (напр., [10]) следует


$|x(t)| \le z(t)$. Из этой оценки получим, что в равенстве


\begin{multline*}


\int_0^t R(t-s) \varphi [ x(s) ] ds=\sum_{\re\lambda_j >


\beta \mu} \bigg\{\int_0^\infty P_j (t-s) e^{\lambda_j (t-s)} \varphi


[ x(s) ] ds - \\


%


-\int_t^\infty P_j (t-s) e^{\lambda_j (t-s)} \varphi [ x(s)


] ds \bigg\} + \sum_{\re\lambda_j < \beta \mu} \int_0^t P_j (t-s)


e^{\lambda_j (t-s)} \varphi [ x(s) ] ds+ \\


%


+\int_0^t R_0(t-s) \varphi [ x(s)] ds


\end{multline*}


три последних интеграла есть $O(t^{\beta (n-1)} e^{\beta \mu t})$,


откуда и вытекает утверждение теоремы.





\begin{cornonum}


Если $\varphi$ --- ограниченная функция, то


\begin{equation}\label{s}


x(t)= \sum_{j=1}^k P_j (t) e^{\lambda_j t} +u(t),


\end{equation}


где $u(t)$ --- ограниченная функция.


\end{cornonum}





Чтобы получить более полную информацию о поведении $x$ при $t \to \infty$,


будем предполагать далее, что


$$


f \in A_m = \bigg\{ y(t)\in C[0,\infty): y(t)=\sum_{j=0}^m \frac{c_j}{(t+1)


^j} + \frac{o(1)}{(t+1)^m} \bigg\},


$$


т.\,е. $f$ имеет тейлоровское разложение на бесконечности порядка $m$.





Рассмотрим сначала неограниченные решения. Из (\ref{s}) следует, что такие


решения либо имеют колебательный характер с неограничено возрастающей


амплитудой, либо имеют вид $x(t)\sim Ct^n e^{\lambda t}$, $\lambda>0$.


Потребуем дополнительно существования


$\lim\limits_{|x| \to \infty} \varphi(x)


=a_0$. Очевидно, можно считать $a_0=0$, т.\,к. слагаемое


$\int\limits_0^t K(t-s)a_0\,ds$


в уравнении (1) можно отнести к свободному члену.





\begin{thm}


Пусть уравнение $1-\wh{K}(z)=0$ имеет в правой полуплоскости $\re z \ge 0$


конечное число корней $\lambda_1,\dotsc,\lambda_k$,


причем $\re\lambda_j > 0$ и


на каждой прямой $\re z=\const$


находится не более двух корней. Пусть $t^m K(t)\in L_1[0,\infty)$,


$t^m K(t)\to 0$ при $t\to\infty$, $f(t)\in A_m$,


$\varphi(x)=\frac{a+o(1)}{\ln^m


(|x|+2)}$ при $|x| \to \infty$. Тогда, если решение уравнения


$(1)$ неограничено,


\begin{equation}


x(t)= \sum_j P_j (t) e^{\lambda_j t} + \sum_{k=0}^m \frac{c_k}{(t+1)^k} +


\frac {o(1)}{(t+1)^m}. \label{x(t)}


\end{equation}


\end{thm}





Действительно, из (\ref{s}) следует, что при достаточно больших $t$


$$


x(t)=P(t)e^{\alpha t}(\sin(\gamma t+\beta(t))+\psi(t)),


$$


где $\alpha>0$, $P(t)>0$, $P(t)\sim ct^k$,


$\beta(t)\to\beta_0$ и $\psi(t)=o(e^{-


\varepsilon_1 t})$ при $t\to\infty$ для некоторого $\varepsilon_1>0$.


Положим $A_1=\{t:|x(t)|\ge P(t)e^{\alpha t-\sqrt{t}}\}$,


$A_2=[0,\infty)\setminus A_1$,


$\varphi_1(t)=\varphi(x(t))\chi_{A_1}(t)+\varphi[P(t)e^{


\alpha t-\sqrt{t}}]\chi_{A_2}(t)$ и


$\varphi_2(t)=\{\varphi(x(t))-\varphi[P(t)e^{\alpha t-\sqrt{t}}


]\}\chi_{A_2}(t)$,


где $\chi_A(t)$ --- характеристическая функция множества $A$.


Можно показать, что ${\int\limits_{A_2}t^m\,dt<\infty}$, и, следовательно,


$t^m\varphi_2(t)\in L_1[0,\infty)$.


Из (1) имеем


$$


x(t)=\int_0^t K(t-s) x(s) ds+f(t)+\int_0^t K(t-s)\varphi_1(s)


ds+\int_0^t \varphi_2(t-s)K(s)ds.


$$


Так как $t^m\varphi_2(t)\in L_1[0,\infty)$, $t^m K(t)\to0$ при $t\to\infty$,


то последний интеграл есть $o(t^{-m})$ при $t\to\infty$ [7]. Далее,


$\varphi_1(t)=\frac{c+o(1)}{t^m}$ и $t^m K(t)\in L_1[0,\infty)$.


Следовательно, второй


интеграл эквивалентен $\frac{c_1}{t^m}$, $t\to\infty$. Таким образом, $x(t)$


является решением линейного уравнения со свободным членом из $A_m$. Отсюда


следует, что $x(t)$ имеет вид (\ref{x(t)}).


Если $x(t)\sim Ce^{\lambda t}$, то представление (\ref{x(t)}) справедливо при


более слабых условиях на $\varphi(x)$.





\begin{thm}


Пусть уравнение $1-\wh{K}(z)=0$ имеет в правой полуплоскости $\re z \ge 0$


конечное число корней $\lambda_j$, причем $\re\lambda_j>0$. Пусть далее


${t^m K(t)\in L_1[0,\infty)}$ и $f(t)\in A_m$.





Eсли $x(t)=e^{\lambda t} ( C+o(1))$, $t \to \infty$,


где $\lambda > 0$, $C \ne 0$, и


$$


\varphi(x)=\sum_{k=1}^m \frac{a_k}{\ln^k ( |x|+2)}+ \frac{o(1)}


{\ln^m ( |x|+2 )},\quad |x| \to \infty,


$$


то $x(t)$ имеет вид $(\ref{x(t)})$.


\end{thm}





Из $x(t)\sim Ce^{\lambda t}$ и (\ref{s}) следует, что


$x(t)=C e^{\lambda t}(1+O(e^{-\varepsilon t}))$, $t\to\infty$, при


некотором $\varepsilon>0$. Тогда $\frac{1}{\ln(|x(t)|+2)}


\in A_m$ и, следовательно, $\varphi[x(t)] \in A_m$. Поэтому


$$


\psi(t)=f(t)+\int_0^t K(t-s)\varphi(x(s))ds\in A_m.


$$


Утверждение теоремы вытекает теперь из того, что $x(t)$ является решением


уравнения


$$


x(t)=\int_0^t K(t-s)x(s)ds+\psi(t).


$$





Перейдем к изучению ограниченных решений. Вообще говоря, нельзя


утверждать, что ограниченное решение имеет предел при $t\to\infty$. Однако,


если $x(t)\to x_0$,


то, как следует из приводимой ниже теоремы, $x(t) \in A_m$.


Подобно тому, как это было сделано с $\varphi(x)$,


здесь также без ограничения


общности можно считать, что $x_0=0$.


Ясно, что в таком случае $\varphi(0)=0$ и


$f(t) \to 0$ при $t \to \infty$.





\begin{thm}


Пусть уравнение $1-\wh{K}(z)=0$ имеет в полуплоскости $\re z\ge0$ конечное


число корней $\lambda_j$, $\re\lambda_j>0$, $t^m K(t) \in L_1[0,\infty)$,


существуют $\varphi^{(m)}(0)$, $\varphi^{\prime \prime}(0)$ и $\varphi^\prime


(0)=0$.


Тогда, если $x(t) \to 0$ при $t \to \infty$, то $x(t) \in A_m$.


\end{thm}





В самом деле, в уравнении (\ref{2}) \ $\psi(t)$


имеет представление (\ref{x(t)}). Отсюда, из (\ref{r}) и (\ref{2}) с


учетом того, что $x(t) \to 0$, найдем


\begin{equation}\label{7}


x(t)= \sum_j \int_t^\infty P_j (t-s) e^{\lambda_j (t-s)}\varphi[x(s)


]ds+\int_0^t R_0(t-s)\varphi[x(s)]ds+


\sum_{j=1}^m \frac{a_j} {(t+1)^j} + \frac{o(1)}{(t+1)^m}.


\end{equation}


Положим $\ov{x}(t)=\sup\limits_{s \ge t} |x(s)|$. Из (\ref{7}), условия


$\varphi(x)=O(x^2)$, $x \to 0$ и монотонности $\ov{x}(t)$ получим


\begin{equation}


\ov{x}(t) \le C_1 \ov{x}^2 (t) +C_2 \ov{x}^2 \bigg(\frac{t}{2} \bigg)+


\frac{C_3} {t+1}.


\end{equation}


Пусть $y(t) = \ov{x}(t)\sqrt{t+1}$. Если $y(t)$ --- неограниченная функция,


то существует такая последовательность $t_n \to \infty$, что $y(t_n)=\sup


\limits_{t \le t_n} y(t) \to \infty$. Тогда


$$


y(t_n) \le \bigg \{ C_1 \ov{x}(t_n) +2 C_2 \ov{x} \bigg (\frac{t_n}{2}


\bigg) \bigg\} y(t_n)+ \frac{C_3} {\sqrt{t_n+1}}.


$$


Так как $\ov{x}(t) \to 0$ при $t \to \infty$, то последнее неравенство


противоречиво при достаточно больших~$n$. Следовательно, $\ov{x}(t)=O(


\frac{1}{\sqrt{t+1}})$, откуда согласно (8) $x(t) = O(\frac{1}


{t+1})$. Так как $t^m R_0(t) \in L_1[0,\infty)$ и $\varphi(x)=O(x^2)$,


$x \to 0$, то первые два слагаемых в (7) являются $O(\frac{1}{(t+1)^2})$


[2] и потому $x(t) \in A_1$. Далее можно показать, что


оператор, определенный правой частью уравнения (7),


переводит функции из $A_k$ в


$A_{k+1}$, $k<m$. Отсюда по индукции получим утверждение теоремы.





Аналогичное утверждение для устойчивого ядра было доказано в [1].





\begin{thebibliography}{99}





\bibitem{1}


Дербен\"ев В.А., Цалюк З.Б. {\em К вопросу об асимптотическом разложении


решений уравнения восстановления} // Дифференц. уравнения. -- 1974. --


Т.\,10. -- \No\,2. -- C.\,335--341. 


\bibitem{2}


Дербен\"ев В.А. {\em Асимптотические свойства решения уравнения


восстановления} // Матем. анализ. -- Краснодар, 1974. -- Вып.\,2.


-- С.\,43--49. 


\bibitem{3}


Дербен\"ев В.А. {\em Асимптотика решения неустойчивого уравнения 


восстановления}. -- Ред. журн. ``Изв. вузов. Математика''.


-- Казань, 1976. -- 13 с. -- Деп. в ВИНИТИ


01.07.76, \No\,2460--76. 


\bibitem{4}


Дербен\"ев В.А., Пуляев В.Ф. {\em Структура резольвенты и устойчивость


линейных интегральных и интегро-дифференциальных уравнений} //


Изв. Сев.-Кавказск. научн. центра высш. школы. Сер. естеств. наук.


-- 1992. -- \No\,1--2. -- С.\,7--14. 


\bibitem{5}


Дербен\"ев В.А., Цалюк З.Б. {\em Асимптотика резольвенты неустойчивого 


уравнения Вольтерра с разностным ядром}


// Матем. заметки. -- 1997. -- Т.\,62. -- \No\,1. -- С.\,88--94. 


\bibitem{6}


Ойнас И.Л., Цалюк З.Б. {\em Асимптотика решений одной системы разностных


уравнений} // Изв. вузов. Математика. -- 2000. -- \No\,4. --  C.\,70--71. 


\bibitem{7}


Цалюк З.Б. {\em О допустимости пары $(Y,X)$ и асимптотике резольвенты для 


системы интегральных уравнений Вольтерра} // Дифференц. уравнения.


-- 1998. -- Т.\,34. -- \No\,9. -- C.\,1226--1230. 


\bibitem{8}


Цалюк З.Б. {\em Асимптотическая структура резольвенты неустойчивого уравнения


Вольтерра с разностным ядром} //


Изв. вузов. Математика. -- 2000. -- \No\,4. -- C.\,50--55. 


\bibitem{9}


Винер Н., Пэли Р. {\em Преобразование Фурье в комплексной области}.


-- М.: Наука, 1964. -- 268 с. 


\bibitem{10}


Азбелев Н.В., Цалюк З.Б. {\em Об интегральных неравенствах}. I //


Матем. сб. -- 1962. -- Т.\,56. -- \No\,3. -- C.\,325--342. 





\end{thebibliography}





\vskip 2cm





\post{\it Кубанский государственный университет }{\it Поступила}


\post{}{10.05.2001}





\label{end}


\end{document}


